\documentclass[a4paper,10pt]{article}

%Packages
% en haut du main.tex
\usepackage{shellesc}
\write18{mkdir -p build}
\pdfoutput=1

\usepackage{txfonts}
\usepackage{titlesec}
\usepackage[utf8]{inputenc}
\usepackage[T1]{fontenc}
\usepackage[french]{babel}
\usepackage{geometry}
\usepackage{graphicx}
\usepackage[intoc]{nomencl}
\usepackage{lipsum}
\usepackage{multicol}
\usepackage{fancyhdr}
\usepackage{hyperref}
\usepackage{float}
\usepackage{setspace}
\usepackage{caption}
\usepackage{booktabs}

% MISE EN PAGE
\geometry{a4paper, top=18mm, bottom=18mm, textwidth=185mm}
\setlength{\columnsep}{6mm} % espace entre les colonnes
\setlength{\parindent}{0pt} % espace indentation

\titleformat{\section}
{\normalfont\bfseries} % <== style du titre
	{\thesection.}{0.5em}{} % numéro + espacement
\titleformat{\subsection}
	{\normalfont\itshape} % <== style du sous-titre
	{\thesubsection.}{0.5em}{} % numéro + espacement

\pagestyle{fancy}
\renewcommand{\footrulewidth}{0pt} %retire la ligne du bas
\fancyhf{}
\rhead{\includegraphics[scale=0.4]{Images/hepia_simple.png}} %ajoute le logo hepia
\rfoot{\textit{\today}} % date en bas à droite
\cfoot{\thepage} %numéro de page

\captionsetup{
  font=footnotesize,   % taille du texte de la légende
  labelsep=colon       % met un ":" après le numéro
}

% NOMENCLATURE
\makenomenclature
\nomenclature{\omega}{la pulsation}
\nomenclature{f}{la fréquence}
\nomenclature{\lambda}{la longueur d'onde}
\nomenclature{k}{le nombre d'onde}
%\nomenclature{$$}{}
% !!! 
%pour enlever la mise en page avec deux colonnes il faut retirer le 
%begin{multiclols}, end{multicols} et les accolades [] qui servent au résumé.
% !!!

\begin{document} %début du document

\begin{multicols*}{2}
[ % l'accolade sert à laisser la mise en page en une seule colonne
% TITRES, AUTEURS, ORGANISATION
%
\begin{center}
\vspace*{0mm}
{\large
% Corriger les fautes dans le titre :
Rapport de laboratoire\\
\vspace{3mm}
Interférences d'ondes acoustiques et de surface   % "accoustiques" → "acoustiques"
}\\
\vspace{8mm}
% Auteur(s) avec leur emails en bas de page
Rozier Vilardell Guilhem\footnote{\textit{\href{guilhem.roziervi@hes-so.ch}{MT2, guilhem.roziervi@hes-so.ch}}},
autre\footnote{\textit{\href{autre@hes-so.ch}{MT2, autre@hes-so.ch}}} 
\\
%
\vspace{3mm}
\footnotesize{
\textit{Haute école de paysagerie, d'ingénierie et d'architecture , Genève, Suisse}
}\\
\vspace{9mm}
\end{center}
%
\rule{\linewidth}{0.4pt} % ligne horizontale
\onehalfspacing
\textbf{Abstract}\\
abstract\\
\noindent\textit{Mots-clés :}mots clés\\
\rule{\linewidth}{0.4pt} % ligne horizontale
]
\singlespacing

% TABLE DES MATIERES
\tableofcontents

% NOMENCLATURE
\printnomenclature

\setlength{\parskip}{1pt} % espace entre deux paragraphes
%Il est mis ici car sinon il change les sauts de la nomenclature et de la talbe des matières

% INTRODUCTION
\section{Introduction}

intro
% THEORIE
\section{Théorie}

théorie


La célérité du son dans un gaz parfait s'écrit :

\begin{equation}
v_{\mathrm{théo}} = \sqrt{\frac{\gamma_{\mathrm{air}}\,R\,T}{M}}
\label{eq:vson}
\end{equation}

avec $\gamma_{\mathrm{air}} = 1{,}4$, $R = 8{,}314$\,J\,mol$^{-1}$\,K$^{-1}$, $M = 0{,}029$\,kg\,mol$^{-1}$ et $T$ la température du milieu, ici $T = 296{,}0$\,K.


% DESCRIPTION EXPERIMENTALE
\section{Description expérimentale et dimensionnement}
%##################GRAPHES##################
\begin{figure}[H]
\centering
%\includegraphics[width=\linewidth]{Images/hepia_simple.png}
\caption{description.}
\label{fig:image}
\end{figure}
%###################EQUATIONS##################
\begin{equation}
    G \approx \frac{95}{55} \approx 1{,}73
\end{equation}

%###################TABLEAUX##################
\begin{table}[H]
\centering
\caption{titre du tableau}
\label{tab:cuve}
\begin{tabular}{ccccc}
\toprule
colonne 1 & colonne 2 & colonne 3 & colonne 4 & colonne 5 \\
\midrule
16{,}76 & 110 & 5 & 12{,}7 & 0{,}213 \\
20{,}29 &  90 & 5 & 10{,}4 & 0{,}211 \\
30{,}00 &  67 & 5 & 7{,}8 & 0{,}232 \\
39{,}80 &  66 & 6 & 6{,}4 & 0{,}253 \\
50{,}30 &  52 & 6 &  5{,}0 & 0{,}251 \\
59{,}90 &  57 & 7 &  4{,}7 & 0{,}281 \\
70{,}30 &  50 & 7 &  4{,}1 & 0{,}290 \\
81{,}10 &  46 & 7 &  4{,}0 & 0{,}308 \\
\bottomrule
\end{tabular}
\end{table}

\section{Résultats et analyse}
résultats et analyse

\section{Conclusion}
CONCLUSION
\end{multicols*}


% ANNEXE
\appendix
\section{Incertitudes}
La loi de propagation des l’incertitudes est donnée par l’équation qui suit :
\begin{equation}
\Delta G = \sqrt{\sum_{i=1}^{N} \left( \frac{\partial G}{\partial x_i}\,\Delta x_i \right)^2}
\end{equation}



\section{Figures et Tables}



FIGURES ET TABLES

\end{document}
